\section{Mirrors and reflective program change}
\label{sec:reflection}

This section refines the schematic transition in
Equation~\eqref{eq:reflective-step}.  It deliberately specifies a semantic
mirror protocol rather than a concrete library API: the Specification requires
mirror-based capability control, transactional change, object-class change,
and activation access, but explicitly leaves the VM-mirror selector set open.
An implementation may choose selectors freely provided that each operation
below is reached only through the corresponding mirror capability.

\subsection{VM cohorts, mirror objects, and authority}

One actor configuration in Equation~\eqref{eq:actor-configuration} is one VM
cohort.  Define its complete object store by
\begin{equation}\label{eq:vm-object-store}
 \mathsf{vmStore}(\Omega)=
 V\uplus\biguplus_{A\in\dom(\mathcal H)}\mathcal H(A).
\end{equation}
The disjointness condition in Equation~\eqref{eq:actor-heap-view} makes this a
map.  Write $\mathsf{ownerActor}_{\Omega}(x)=A$ exactly when
$x\in\dom(\mathcal H(A))$, and write
$\mathsf{ownerActor}_{\Omega}(x)=\kw{shared}$ exactly when $x\in\dom(V)$.
These two cases are disjoint.

A mirror target and a reflective right are drawn from
\begin{equation}\label{eq:mirror-targets-rights}
\begin{split}
 \zeta::={}&\kw{vmTarget}(w)\mid\kw{programTarget}
  \mid\kw{mixinTarget}(M)\mid\kw{methodTarget}(i_f)\\
 &\mid\kw{classTarget}(C)\mid\kw{objectTarget}(o)
  \mid\kw{activationTarget}(a)\mid\kw{actorTarget}(A),\\
 \varrho\in\mathit{ReflectRight}={}&
 \{\kw{inspect},\kw{editCode},\kw{editClass},\kw{reclassifyObject},\\
 &\quad\kw{editObjectSlots},\kw{editActivation},\kw{controlContinuation},\\
 &\quad\kw{controlActorStack},\\
 &\quad
       \kw{primitive},\kw{delegate}\}.
\end{split}
\end{equation}
Extend the heap-record grammar of Equation~\eqref{eq:heap-records} by
\begin{equation}\label{eq:mirror-record}
 H(g)=\kw{mirror}\langle C_g,w,\zeta,R\rangle,
 \qquad R\in\mathcal P_{\mathrm{fin}}(\mathit{ReflectRight}),
\end{equation}
where $C_g$ is a platform mirror class.  Thus $\classof_H(g)=C_g$.
Mirror identities are fresh and unforgeable.  A mirror is actor-local mutable
state, never a value; ordinary cross-actor transfer therefore yields a far
reference under Rules \rulelabel{RR-Near-Far} and \rulelabel{RR-Far-Pass}.
Possession of such a reference can still confer authority, because the
eventual request executes in the mirror's owning actor.

Target containment is the least reflexive relation $\sqsubseteq_{P,H,w}$
satisfying
\begin{equation}\label{eq:mirror-target-containment}
\begin{aligned}
 \zeta&\sqsubseteq_{P,H,w}\kw{vmTarget}(w),\\
 \kw{mixinTarget}(M)&\sqsubseteq_{P,H,w}\kw{programTarget},\\
 \kw{methodTarget}(i_f)&\sqsubseteq_{P,H,w}\kw{mixinTarget}(M)
   &&\text{if }\mathsf{methodOwner}_P(i_f)=M,\\
 \kw{classTarget}(C)&\sqsubseteq_{P,H,w}
       \kw{mixinTarget}(\mixin_H(C)),\\
 \kw{objectTarget}(o)&\sqsubseteq_{P,H,w}
       \kw{classTarget}(\classof_H(o)),\\
 \kw{activationTarget}(a)&\sqsubseteq_{P,H,w}
       \kw{objectTarget}(H(a).o).
\end{aligned}
\end{equation}
The first, generic clause places an actor target occurring in this cohort below
its VM target.  An actor target does not contain the actor's whole heap:
activation and object mirrors remain separate capabilities.  It authorizes
only the stack operations explicitly requiring \kw{controlActorStack} below.
The method owner used here is the partial function
\begin{equation}\label{eq:method-owner}
 \mathsf{methodOwner}_P(i_f)=M
 \quad\Longleftrightarrow\quad
 \exists s,f.\;\methods_P(M)(s)=f\land\mathsf{methodId}(f)=i_f.
\end{equation}
Uniqueness follows from stable method identities.  VM containment includes all
targets in Equation~\eqref{eq:mirror-targets-rights} that occur in the store of
that VM; it does not include any target belonging to another VM.

The platform may create an initial VM mirror only as part of the authority
explicitly supplied to an application.  There is no operation that obtains it
from an ambient global.  A holder may attenuate, but not amplify, authority:
\[
\dfrac{\begin{gathered}
 H_A(g)=\kw{mirror}\langle C_g,w,\zeta,R\rangle\qquad
 \kw{delegate}\in R\\
 \zeta'\sqsubseteq_{P,\widehat H_A,w}\zeta\qquad R'\subseteq R
 \qquad g'\text{ fresh}\\
 H_A'=H_A[g'\mapsto\kw{mirror}\langle C_g,w,\zeta',R'\rangle]
\end{gathered}}
 {\Omega,A,g\vdash\mathsf{deriveMirror}(\zeta',R')
       \Downarrow\langle\Omega[\mathcal H(A)\mapsto H_A'],g'\rangle}
\quad(\rulelabel{Mirror-Derive})
\]
Access modifiers play no role in this rule.  A private method is visible to a
method or enclosing-mixin mirror precisely because the mirror, rather than a
source-level send, is the authority check.

\subsection{Reified control and stack images}
\label{sec:debugging}

The sequential notation of Equation~\eqref{eq:configuration} stores the active
term once, beside the activation stack.  Debugging requires a uniform
per-frame view.  Let $b_c\in\mathit{CodeImageId}$ name an immutable source,
bytecode, or native-code image and define a control location by
\begin{equation}\label{eq:control-location-domain}
 \lambda_c::=\kw{sourcePoint}(z)\mid
   \kw{machinePoint}(b_c,n_c),\qquad n_c\in\mathbb N.
\end{equation}
Introduce the internal waiting term \kw{awaitResult} and the control-image
domain
\begin{equation}\label{eq:frame-control-domain}
 \kappa::=\kw{activeControl}\langle\pi,t\rangle
       \mid\kw{suspendedControl}\langle\pi\rangle.
\end{equation}
The current term and source-level control location of a frame are the total
functions
\begin{equation}\label{eq:frame-current-term-and-location}
\begin{aligned}
 \mathsf{currentTerm}(\kw{activeControl}\langle\pi,t\rangle)&=t,\\
 \mathsf{controlLocation}(\kw{activeControl}\langle\pi,t\rangle)
   &=\mathsf{site}(\pi,t),\\
 \mathsf{currentTerm}(\kw{suspendedControl}\langle\pi\rangle)
   &=\kw{awaitResult},\\
 \mathsf{controlLocation}(\kw{suspendedControl}\langle\pi\rangle)
   &=\mathsf{resumeSite}(\pi).
\end{aligned}
\end{equation}
$\mathsf{site}$ and $\mathsf{resumeSite}$ are the control-map functions emitted
by elaboration or compilation and return a $\lambda_c$.  A source evaluator
normally returns \kw{sourcePoint}; a bytecode or native implementation returns
\kw{machinePoint}, whose second component is its actual program counter.
Materialized terms and frames are not changed by a later code
install, so the map for a live or debugger-retained frame must be retained as
long as that frame is resumable.

For
$\mathcal A=\langle a_1,\pi_1\rangle\cdots
 \langle a_m,\pi_m\rangle$ with $m>0$, define its canonical control image by
\begin{equation}\label{eq:stack-control-image}
\begin{split}
 \mathsf{controlImage}(\mathcal A,t)=
 &\langle a_1,\kw{suspendedControl}\langle\pi_1\rangle\rangle\cdots\\[-1mm]
 &\langle a_{m-1},\kw{suspendedControl}\langle\pi_{m-1}\rangle\rangle
 \cdot\langle a_m,\kw{activeControl}\langle\pi_m,t\rangle\rangle.
\end{split}
\end{equation}
Thus every live frame has a current term: a lower frame is explicitly waiting
for its callee's result at the location determined by its continuation stack.
The inverse partial function is
\begin{equation}\label{eq:decode-control-image}
 \mathsf{decodeControlImage}
   (\mathsf{controlImage}(\mathcal A,t))=\langle\mathcal A,t\rangle.
\end{equation}
It is undefined for an empty image, for an active frame below the top, or for a
suspended top frame.

A reified frame contains both the activation record and its control image:
\begin{equation}\label{eq:debug-frame-image}
 \varphi=\kw{frameImage}\langle a,H(a),\kappa\rangle.
\end{equation}
The actor-stack snapshot is the immutable value
\begin{equation}\label{eq:actor-stack-snapshot}
\begin{split}
 \mathsf{actorStackSnapshot}_{\Omega}(B)
   &\simeq\langle\kw{running},
      \mathsf{reifyFrames}_{U}(
        \mathsf{controlImage}(\mathcal A,t))\rangle\\[-1mm]
   &\hspace{30mm}\text{if }\mathcal X(B)=
      \kw{runningTurn}\langle\mathcal A,t,p,u\rangle,\\
 \mathsf{actorStackSnapshot}_{\Omega}(B)
   &\simeq\langle\kw{paused}(r_s,d_s),
      \mathsf{reifyFrames}_{U}(
        \mathsf{controlImage}(\mathcal A,t))\rangle\\[-1mm]
   &\hspace{30mm}\text{if }\mathcal X(B)=
      \kw{pausedTurn}\langle\mathcal A,t,p,u,r_s,d_s\rangle,
\end{split}
\end{equation}
where $U=\mathsf{vmStore}(\Omega)$ and $\mathsf{reifyFrames}$ replaces every
$\langle a,\kappa\rangle$ by Equation~\eqref{eq:debug-frame-image}.  It is
undefined for an idle actor.  The paused form exposes the unforgeable token
that must accompany any later paused-stack mutation.

The unique live control of an activation is
\begin{equation}\label{eq:live-activation-control}
\begin{split}
 \mathsf{liveControl}_{\Omega}(a)=\kappa
 \quad\Longleftrightarrow\quad&
 \exists B,\bar\varphi_0,\bar\varphi_1,R.\\[-1mm]
 &\mathsf{actorStackSnapshot}_{\Omega}(B)=
 \langle m,\bar\varphi_0\cdot
   \langle\kw{frameImage}\langle a,R,\kappa\rangle\rangle
   \cdot\bar\varphi_1\rangle.
\end{split}
\end{equation}
Uniqueness follows from actor ownership and the prohibition on duplicate live
activation identities.  The function is undefined for a retained activation
that is no longer on a stack.

\paragraph{Editable stack templates.}
A debugger constructs an immutable template before requesting an atomic
install.  Template-local keys and frames are
\begin{equation}\label{eq:stack-template-domain}
 \eta::=\kw{retain}(a)\mid\kw{freshFrame}(u_f),
 \qquad
 \psi=\kw{frameTemplate}\langle\eta,R,\kappa\rangle,
 \qquad
 \mathcal S=\langle\psi_1,\ldots,\psi_m\rangle.
\end{equation}
$R$ is a complete activation-record template.  Its continuation, home, and
lexical-environment activation references may name other keys in the same
$\mathcal S$.  A retained key preserves an existing activation identity; a
fresh-frame key allocates a new identity at installation.  Consequently a
copied frame is never made live twice under one activation identity.

The three basic stack-editing operations have distinct names.  First, promote
and demote control without overloading either operation:
\begin{equation}\label{eq:promote-demote-frame-control}
\begin{aligned}
 &\mathsf{promoteFrame}(\kw{frameTemplate}\langle
      \eta,R,\kw{suspendedControl}\langle\pi\rangle\rangle,v)
   \\[-1mm]&\qquad={}\kw{frameTemplate}\langle
      \eta,R,\kw{activeControl}\langle\pi,v\rangle\rangle,\\
 &\mathsf{demoteFrame}(\kw{frameTemplate}\langle
      \eta,R,\kw{activeControl}\langle\pi,t\rangle\rangle)
   \\[-1mm]&\qquad={}\kw{frameTemplate}\langle
      \eta,R,\kw{suspendedControl}\langle\pi\rangle\rangle.
\end{aligned}
\end{equation}
Popping $j$ frames supplies the chosen continuation value to the new top:
\begin{equation}\label{eq:pop-stack-frames}
 \mathsf{popStackFrames}
   (\mathcal S_0\cdot\langle\psi\rangle\cdot\mathcal S_1,
      |\mathcal S_1|,v)
 =\mathcal S_0\cdot\langle\mathsf{promoteFrame}(\psi,v)\rangle,
 \qquad |\mathcal S_1|<|\mathcal S|.
\end{equation}
Pushing a prepared active frame makes the old top wait for the pushed frame's
result:
\begin{equation}\label{eq:push-stack-frame}
 \mathsf{pushStackFrame}
  (\mathcal S_0\cdot\langle\psi\rangle,\psi')
 =\mathcal S_0\cdot\langle\mathsf{demoteFrame}(\psi),\psi'\rangle,
 \quad
 \mathsf{currentTerm}(\psi'.\kappa)\ne\kw{awaitResult}.
\end{equation}
Finally, copying a contiguous segment allocates fresh template keys and rewires
every stack-structural activation link internal to the copied segment:
\begin{equation}\label{eq:copy-stack-frame-segment}
\begin{aligned}
 \mathsf{copyStackFrameSegment}(\mathcal S,i,j,\bar u)
   &=\mathsf{rename}_{\nu}
       (\langle\psi_i,\ldots,\psi_j\rangle),\\
 \nu(\psi_k.\eta)&=\kw{freshFrame}(u_k)\qquad(i\le k\le j),
 \qquad |\bar u|=j-i+1.
\end{aligned}
\end{equation}
Continuation, home, and lexical-environment links outside the copied segment
remain references to their original retained activations.  Parameter and local
values are shallow-copied; an activation object stored there is not implicitly
retargeted.  A debugger may concatenate the returned segment at any position,
edit its records or controls, and submit the complete result for validation.
These pure operations therefore express push, pop, single-frame copy, and
multi-frame copy without making a partially edited live stack observable.

\subsection{Inspection}

Mirror queries are
\begin{equation}\label{eq:mirror-query-grammar}
\begin{split}
 q_r::={}&\kw{programVersionQuery}\mid\kw{mixinDescriptorQuery}(M)\\
 &\mid\kw{methodDescriptorQuery}(i_f)\mid\kw{classDescriptorQuery}(C)\\
 &\mid\kw{objectDescriptorQuery}(o)\mid\kw{activationDescriptorQuery}(a)\\
 &\mid\kw{activationControlQuery}(a)\mid\kw{currentActivationQuery}(A)\\
 &\mid\kw{actorStackQuery}(A)\mid\kw{allInstancesOfQuery}(C)\\
 &\mid\kw{mixinApplicationsQuery}(M).
\end{split}
\end{equation}
For a running actor, the current activation is
\begin{equation}\label{eq:current-activation}
 \mathsf{currentActivation}_{\Omega}(A)=a
 \quad\Longleftrightarrow\quad
 \begin{aligned}[t]
 \mathcal X(A)&=\kw{runningTurn}\langle
   \mathcal A_0\mathbin{\cdot}\langle a,\pi\rangle,t,p,u\rangle
 \quad\lor\\
 \mathcal X(A)&=\kw{pausedTurn}\langle
   \mathcal A_0\mathbin{\cdot}\langle a,\pi\rangle,t,p,u,r_s,d_s\rangle.
 \end{aligned}
\end{equation}
It is undefined for an idle actor.  Descriptor construction is given by the
following separate functions, each of which returns an immutable value object:
\begin{equation}\label{eq:reflective-descriptors}
\begin{aligned}
 \mathsf{describeMixin}_{P}(M)&\simeq
   \left\langle\begin{gathered}\decl_P(M),\methods_P(M),
       \mathsf{ownSlots}_P(M),\\
       P.\mathit{mixin}(M).\mathit{nestedDecls},\iota_P(M)
   \end{gathered}\right\rangle,\\
 \mathsf{describeMethod}_{P}(i_f)
   &\simeq\langle M,f,P.\mathit{body}(i_f)\rangle,\\[-1mm]
 &\quad\text{where }M=\mathsf{methodOwner}_P(i_f),\quad
       \methods_P(M)(\selector(f))=f,\\
 \mathsf{describeClass}_{H}(C)&\simeq
   \left\langle\begin{gathered}\mixin_H(C),\super_H(C),H(C).eo,\\
       H(C).C_m,H(C).q,H(C).\chi_C
   \end{gathered}\right\rangle,\\
 \mathsf{describeObject}_{H}(o)
   &\simeq\langle\classof_H(o),H(o).\sigma,H(o).\eta\rangle,\\
 \mathsf{describeActivation}_{H}(a)&\simeq
   \left\langle\begin{gathered}H(a).p,H(a).\rho,H(a).\ell,H(a).o,H(a).C,\\
       H(a).k,H(a).h,H(a).q_a,H(a).\gamma
   \end{gathered}\right\rangle.
\end{aligned}
\end{equation}
The complete debugger-facing activation descriptor adds its current control
when the activation is live:
\begin{equation}\label{eq:debug-activation-descriptor}
 \mathsf{describeDebugActivation}_{\Omega}(a)
  \simeq\langle
    \mathsf{describeActivation}_{\mathsf{vmStore}(\Omega)}(a),
    \mathsf{option}(\mathsf{liveControl}_{\Omega}(a))\rangle.
\end{equation}
$\mathsf{option}$ is \kw{none} when
Equation~\eqref{eq:live-activation-control} is undefined.  Thus retained dead
activations remain inspectable, while every live activation exposes both its
current term and the source/PC location of
Equation~\eqref{eq:frame-current-term-and-location}.
The symbol $\simeq$ means an immutable, recursively reified snapshot, not a
mutable alias to a metalevel record.  In particular, a continuation component
is the activation identity specified in Section~3, reified as an activation
mirror only if the caller has authority for that target.

Fix an implementation-chosen total order $\prec$ on the object identities of
one VM cohort.  The exact-instance query returns a snapshot sequence, not a
live view:
\begin{equation}\label{eq:all-instances-of}
 \mathsf{allInstancesOf}_{\Omega}(C)=
 \mathsf{sort}_{\prec}
 \{x\in\dom(\mathsf{vmStore}(\Omega))
       \mid \classof_{\mathsf{vmStore}(\Omega)}(x)=C\}.
\end{equation}
Thus subclasses are excluded.  Presence in the VM store is the abstract
notion of liveness.  The optional garbage-collection layer makes this precise
by removing eligible identities from that store; without that layer, the query
observes every exact instance allocated so far.  The order is deliberately a
platform choice because the mirror protocol promises no allocation order.

The applications of a mixin are the live class objects whose class records
name that exact static mixin identity:
\begin{equation}\label{eq:mixin-applications}
 \mathsf{mixinApplications}_{\Omega}(M)=
 \mathsf{sort}_{\prec}
 \{C\in\dom(\mathsf{vmStore}(\Omega))
       \mid \mathsf{isClass}_{\mathsf{vmStore}(\Omega)}(C)
       \land \mixin_{\mathsf{vmStore}(\Omega)}(C)=M\}.
\end{equation}
This is an exact-mixin query: applications of other mixins and ordinary
instances of the resulting classes are excluded.  A metaclass is included
when its class record applies $M$, so a class-side mixin is treated uniformly
with an instance-side mixin.  As with Equation~\eqref{eq:all-instances-of}, the
result is a snapshot in the fixed order $\prec$.  Working Specification 0.109
requires both population queries.

For $\Omega=\langle P,V,\mathcal H,\mathcal X,\mathcal Q,\mathcal N,
\Pi,\Phi,\Gamma,\mathcal E,\mathcal D,w,n\rangle$, put
$U=\mathsf{vmStore}(\Omega)$.  The authority target of every query is
\begin{equation}\label{eq:mirror-query-target}
\begin{aligned}
 \mathsf{queryTarget}_{P,U,w}(\kw{programVersionQuery})
   &=\kw{vmTarget}(w),\\
 \mathsf{queryTarget}_{P,U,w}(\kw{mixinDescriptorQuery}(M))
   &=\kw{mixinTarget}(M),\\
 \mathsf{queryTarget}_{P,U,w}(\kw{methodDescriptorQuery}(i_f))
   &=\kw{methodTarget}(i_f),\\
 \mathsf{queryTarget}_{P,U,w}(\kw{classDescriptorQuery}(C))
   &=\kw{classTarget}(C),\\
 \mathsf{queryTarget}_{P,U,w}(\kw{objectDescriptorQuery}(o))
   &=\kw{objectTarget}(o),\\
 \mathsf{queryTarget}_{P,U,w}(\kw{activationDescriptorQuery}(a))
   &=\kw{activationTarget}(a),\\
 \mathsf{queryTarget}_{P,U,w}(\kw{activationControlQuery}(a))
   &=\kw{activationTarget}(a),\\
 \mathsf{queryTarget}_{P,U,w}(\kw{currentActivationQuery}(A))
   &=\kw{actorTarget}(A),\\
 \mathsf{queryTarget}_{P,U,w}(\kw{actorStackQuery}(A))
   &=\kw{actorTarget}(A),\\
 \mathsf{queryTarget}_{P,U,w}(\kw{allInstancesOfQuery}(C))
   &=\kw{classTarget}(C),\\
 \mathsf{queryTarget}_{P,U,w}(\kw{mixinApplicationsQuery}(M))
   &=\kw{mixinTarget}(M).
\end{aligned}
\end{equation}
Its result is the corresponding immutable snapshot:
\begin{equation}\label{eq:mirror-query-result}
\begin{aligned}
 \mathsf{queryResult}_{\Omega}(\kw{programVersionQuery})&=n,\\
 \mathsf{queryResult}_{\Omega}(\kw{mixinDescriptorQuery}(M))
   &=\mathsf{describeMixin}_{P}(M),\\
 \mathsf{queryResult}_{\Omega}(\kw{methodDescriptorQuery}(i_f))
   &=\mathsf{describeMethod}_{P}(i_f),\\
 \mathsf{queryResult}_{\Omega}(\kw{classDescriptorQuery}(C))
   &=\mathsf{describeClass}_{U}(C),\\
 \mathsf{queryResult}_{\Omega}(\kw{objectDescriptorQuery}(o))
   &=\mathsf{describeObject}_{U}(o),\\
 \mathsf{queryResult}_{\Omega}(\kw{activationDescriptorQuery}(a))
   &=\mathsf{describeDebugActivation}_{\Omega}(a),\\
 \mathsf{queryResult}_{\Omega}(\kw{activationControlQuery}(a))
   &=\mathsf{liveControl}_{\Omega}(a),\\
 \mathsf{queryResult}_{\Omega}(\kw{currentActivationQuery}(A))
   &=\mathsf{currentActivation}_{\Omega}(A),\\
 \mathsf{queryResult}_{\Omega}(\kw{actorStackQuery}(A))
   &=\mathsf{actorStackSnapshot}_{\Omega}(A),\\
 \mathsf{queryResult}_{\Omega}(\kw{allInstancesOfQuery}(C))
   &=\mathsf{allInstancesOf}_{\Omega}(C),\\
 \mathsf{queryResult}_{\Omega}(\kw{mixinApplicationsQuery}(M))
   &=\mathsf{mixinApplications}_{\Omega}(M).
\end{aligned}
\end{equation}
Inspection is the partial judgment
\[
\dfrac{\begin{gathered}
 \widehat H_A(g)=\kw{mirror}\langle C_g,w,\zeta,R\rangle\qquad
 \kw{inspect}\in R\\
 \mathsf{queryTarget}_{P,U,w}(q_r)
       \sqsubseteq_{P,U,w}\zeta\qquad
 d=\mathsf{queryResult}_{\Omega}(q_r)
\end{gathered}}
 {\Omega,A,g\vdash\mathsf{inspectMirror}(q_r)\Downarrow d}
\quad(\rulelabel{Mirror-Inspect})
\]
There is no failed-premise fallback that reveals a target.  The concrete
mirror library may reify denial as an exception, but must not disclose the
descriptor.

\subsection{Reflective transactions}

A transaction $\Delta$ is a finite sequence of commands, each explicitly
carrying the mirror that authorizes it:
\begin{equation}\label{eq:debug-stop-reasons}
 d_s::=\kw{exceptionStop}(x)\mid\kw{breakpointStop}(\lambda_c)
       \mid\kw{requestedStop}(o).
\end{equation}
\begin{equation}\label{eq:reflection-command-grammar}
\begin{split}
 \delta_{\mathrm{code}}::={}&
 \kw{replaceMethodBody}(g,i_f,b)
 \mid\kw{replaceMethodDefinition}(g,M,i_f,f)\\
 &\mid\kw{addMethodDefinition}(g,M,f)
 \mid\kw{removeMethodDefinition}(g,M,i_f)\\
 &\mid\kw{replaceSlotDeclarations}(g,M,\delta_s)
 \mid\kw{replaceNestedDeclarations}(g,M,\delta_n)\\
 &\mid\kw{replaceMixinInitializer}(g,M,\iota),\\
 \delta_{\mathrm{class}}::={}&
 \kw{changeSuperclass}(g,C,S)
 \mid\kw{changeClassEnclosingObject}(g,C,eo),\\
 \delta_{\mathrm{objectClass}}::={}&
 \kw{changeObjectClass}(g,o,C),\\
 \delta_{\mathrm{objectSlot}}::={}&
 \kw{writeObjectSlot}(g,o,j,v),\\
 \delta_{\mathrm{activation}}::={}&
 \kw{writeActivationParameter}(g,a,x,v)
 \mid\kw{writeActivationLocal}(g,a,j,v)\\
 &\mid\kw{changeActivationCurrentClass}(g,a,C)\\
 &\mid\kw{changeActivationContinuation}(g,a,k)
 \mid\kw{markActivationUncontinuable}(g,a),\\
 \delta_{\mathrm{continuation}}::={}&
 \kw{continueAtActivation}(g,A,a,v),\\
 \delta_{\mathrm{debug}}::={}&
 \kw{pauseActor}(g,B,d_s)\\
 &\mid\kw{replaceCurrentActorStack}(g,B,\mathcal S)\\
 &\mid\kw{replacePausedActorStack}(g,B,r_s,\mathcal S)\\
 &\mid\kw{resumePausedActorAtFullSpeed}(g,B,r_s)\\
 &\mid\kw{replaceAndResumePausedActorStack}(g,B,r_s,\mathcal S),\\
 \delta_r::={}&\delta_{\mathrm{code}}
 \mid\delta_{\mathrm{class}}
 \mid\delta_{\mathrm{objectClass}}
 \mid\delta_{\mathrm{objectSlot}}
 \mid\delta_{\mathrm{activation}}
 \mid\delta_{\mathrm{continuation}}
 \mid\delta_{\mathrm{debug}},\\
 \Delta::={}&\langle\delta_1,\ldots,\delta_m\rangle.
\end{split}
\end{equation}
By definition, $\mathit{CodeCommand}$ is generated by
$\delta_{\mathrm{code}}$, and $\mathit{ClassCommand}$ by
$\delta_{\mathrm{class}}$.  The sets $\mathit{ObjectClassCommand}$ and
$\mathit{ObjectSlotCommand}$ are generated by
$\delta_{\mathrm{objectClass}}$ and $\delta_{\mathrm{objectSlot}}$,
respectively.  Finally, $\mathit{ActivationCommand}$ is generated by
$\delta_{\mathrm{activation}}$, and $\mathit{ContinuationCommand}$ by
$\delta_{\mathrm{continuation}}$, and $\mathit{DebuggerCommand}$ by
$\delta_{\mathrm{debug}}$.  These seven sets are disjoint by outer
constructor and form an exhaustive partition of $\delta_r$.  In particular,
a class command edits a runtime class record, an
object-class command reclassifies one object, and an object-slot command
changes state within that object's existing final layout.  These are distinct
operations with distinct rights and distinct preparation phases.  The names
are intentionally not overloaded: changing a whole method definition is also
distinct from replacing only its body.

Define $\mathsf{commandMirror}(\delta_r)$ as the first field of every command.
The required right is the total function
\begin{equation}\label{eq:reflection-command-right}
 \mathsf{requiredRight}(\delta_r)=
 \begin{cases}
  \kw{editCode}&\delta_r\in\mathit{CodeCommand},\\
  \kw{editClass}&\delta_r\in\mathit{ClassCommand},\\
  \kw{reclassifyObject}&\delta_r\in\mathit{ObjectClassCommand},\\
  \kw{editObjectSlots}&\delta_r\in\mathit{ObjectSlotCommand},\\
  \kw{editActivation}&\delta_r\in\mathit{ActivationCommand},\\
  \kw{controlContinuation}&\delta_r\in\mathit{ContinuationCommand},\\
  \kw{controlActorStack}&\delta_r\in\mathit{DebuggerCommand}.
 \end{cases}
\end{equation}
The target function is
\begin{equation}\label{eq:reflection-command-target}
 \mathsf{commandTarget}(\delta_r)=
 \begin{cases}
  \kw{methodTarget}(i_f)&\delta_r\text{ names an existing }i_f,\\
  \kw{mixinTarget}(M)&\delta_r\text{ adds a method or edits mixin structure},\\
  \kw{classTarget}(C)&\delta_r\in\mathit{ClassCommand},\\
  \kw{objectTarget}(o)&\delta_r\in
       \mathit{ObjectClassCommand}\cup\mathit{ObjectSlotCommand},\\
  \kw{activationTarget}(a)&\delta_r\in
       \mathit{ActivationCommand}\cup\mathit{ContinuationCommand},\\
  \kw{actorTarget}(B)&\delta_r\in\mathit{DebuggerCommand}
       \text{ and names actor }B.
 \end{cases}
\end{equation}
For method addition the target is its mixin; for removal, definition
replacement, and body replacement the method identity is the target.  The
authorization predicate is
\begin{equation}\label{eq:reflection-authorized}
\begin{split}
 \mathsf{authorized}_{\Omega,A}(\delta_r)\quad\Longleftrightarrow\quad{}
 &\widehat H_A(g)=\kw{mirror}\langle C_g,w,\zeta,R\rangle\ \land\\
 &g=\mathsf{commandMirror}(\delta_r)\ \land\
 \mathsf{requiredRight}(\delta_r)\in R\ \land\\
 &\mathsf{commandTarget}(\delta_r)
   \sqsubseteq_{P,\mathsf{vmStore}(\Omega),w}\zeta.
\end{split}
\end{equation}
Thus syntactic possession of a mirror variable is not sufficient: the
unforgeable heap record, VM identity, target containment, and operation right
must all agree.

Transactions are simultaneous.  Let $\lambda$ range over the tagged write
locations below, and define the complete write-footprint function
$\mathsf{writes}$ by
{\footnotesize
\begin{equation}\label{eq:stack-template-record-locations}
 \mathsf{templateRecordLocations}(\mathcal S)=
 \{\kw{activationRecordLoc}(a)\mid
   \kw{frameTemplate}\langle\kw{retain}(a),R,\kappa\rangle
      \text{ occurs in }\mathcal S\}.
\end{equation}
}
{\small
\begin{equation}\label{eq:reflection-command-footprint}
\begin{aligned}
 &\mathsf{writes}(\kw{replaceMethodBody}(g,i_f,b))\\
 &\qquad=\{\kw{methodBodyLoc}(i_f)\},\\
 &\mathsf{writes}(\kw{replaceMethodDefinition}(g,M,i_f,f))\\
 &\qquad=\{\kw{methodDefLoc}(i_f),\kw{methodBodyLoc}(i_f),
       \kw{selectorLoc}(M,\selector(f))\},\\
 &\mathsf{writes}(\kw{addMethodDefinition}(g,M,f))\\
 &\qquad=\{\kw{methodDefLoc}(\mathsf{methodId}(f)),
       \kw{selectorLoc}(M,\selector(f))\},\\
 &\mathsf{writes}(\kw{removeMethodDefinition}(g,M,i_f))\\
 &\qquad=\{\kw{methodDefLoc}(i_f),\kw{methodBodyLoc}(i_f)\},\\
 &\mathsf{writes}(\kw{replaceSlotDeclarations}(g,M,\delta_s))
  =\{\kw{slotDeclsLoc}(M)\},\\
 &\mathsf{writes}(\kw{replaceNestedDeclarations}(g,M,\delta_n))
  =\{\kw{nestedDeclsLoc}(M)\},\\
 &\mathsf{writes}(\kw{replaceMixinInitializer}(g,M,\iota))
  =\{\kw{mixinInitializerLoc}(M)\},\\
 &\mathsf{writes}(\kw{changeSuperclass}(g,C,S))
  =\{\kw{superclassLoc}(C)\},\\
 &\mathsf{writes}(\kw{changeClassEnclosingObject}(g,C,eo))
  =\{\kw{classEnclosingLoc}(C)\},\\
 &\mathsf{writes}(\kw{changeObjectClass}(g,o,C))
  =\{\kw{objectClassLoc}(o)\},\\
 &\mathsf{writes}(\kw{writeObjectSlot}(g,o,j,v))
  =\{\kw{objectSlotLoc}(o,j)\},\\
 &\mathsf{writes}(\kw{writeActivationParameter}(g,a,x,v))\\
 &\qquad=\{\kw{activationParameterLoc}(a,x)\},\\
 &\mathsf{writes}(\kw{writeActivationLocal}(g,a,j,v))
  =\{\kw{activationLocalLoc}(a,j)\},\\
 &\mathsf{writes}(\kw{changeActivationCurrentClass}(g,a,C))\\
 &\qquad=\{\kw{activationClassLoc}(a)\},\\
 &\mathsf{writes}(\kw{changeActivationContinuation}(g,a,k))\\
 &\qquad=\{\kw{activationContinuationLoc}(a)\},\\
 &\mathsf{writes}(\kw{markActivationUncontinuable}(g,a))\\
 &\qquad=\{\kw{activationContinuableLoc}(a)\},\\
 &\mathsf{writes}(\kw{continueAtActivation}(g,A,a,v))\\
 &\qquad=\{\kw{actorStackLoc}(A),\kw{activationContinuationLoc}(a)\},\\
 &\mathsf{writes}(\kw{pauseActor}(g,B,d_s))
   =\{\kw{actorStackLoc}(B)\},\\
 &\mathsf{writes}(\kw{replaceCurrentActorStack}(g,B,\mathcal S))\\
 &\qquad=\{\kw{actorStackLoc}(B)\}\cup
      \mathsf{templateRecordLocations}(\mathcal S),\\
 &\mathsf{writes}(\kw{replacePausedActorStack}(g,B,r_s,\mathcal S))\\
 &\qquad=\{\kw{actorStackLoc}(B)\}\cup
      \mathsf{templateRecordLocations}(\mathcal S),\\
 &\mathsf{writes}(\kw{resumePausedActorAtFullSpeed}(g,B,r_s))
   =\{\kw{actorStackLoc}(B)\},\\
 &\mathsf{writes}(\kw{replaceAndResumePausedActorStack}
        (g,B,r_s,\mathcal S))\\
 &\qquad=\{\kw{actorStackLoc}(B)\}\cup
      \mathsf{templateRecordLocations}(\mathcal S).
\end{aligned}
\end{equation}
}
Location overlap is the least reflexive, symmetric relation
\begin{equation}\label{eq:reflection-location-overlap}
\begin{aligned}
 \lambda&\mathrel{\bowtie}\lambda,\\
 \kw{activationRecordLoc}(a)&\mathrel{\bowtie}
   \kw{activationFieldLoc}(a,\xi),\\
 \kw{activationFieldLoc}(a,\xi)&\mathrel{\bowtie}
   \kw{activationRecordLoc}(a).
\end{aligned}
\end{equation}
Each concrete parameter, local, class, continuation, and continuability
location above is the corresponding instance of
$\kw{activationFieldLoc}(a,\xi)$.  All other distinct tagged locations are
nonoverlapping.  This lets two edits
to different fields of one activation coexist while making a whole-record
stack installation conflict with either edit.
For $\Delta=\langle\delta_1,\ldots,\delta_m\rangle$, conflict and
nonconflict are therefore the predicates
\begin{equation}\label{eq:reflection-conflict}
\begin{aligned}
 \mathsf{conflict}_{\Delta}(i,j)
  &\Longleftrightarrow 1\le i<j\le m\ \land\
  &\qquad\exists\lambda\in\mathsf{writes}(\delta_i),
          \lambda'\in\mathsf{writes}(\delta_j).
          \ \lambda\mathrel{\bowtie}\lambda',\\
 \mathsf{nonconflicting}(\Delta)
  &\Longleftrightarrow
    \neg\exists i,j.\;\mathsf{conflict}_{\Delta}(i,j).
\end{aligned}
\end{equation}
The multi-location cases make whole-method replacement overlap body
replacement, and make continuation transfer overlap another transfer for the
same actor or a direct mutation of the resumed continuation.

\subsection{Source change and re-elaboration}

The code-command projection is
\begin{equation}\label{eq:code-command-projection}
 \mathsf{codeCommands}(\Delta)=
 \mathsf{filter}_{\mathit{CodeCommand}}(\Delta).
\end{equation}
Its direct source effect is the partial function $\mathsf{patchSource}$:
\begin{equation}\label{eq:source-patch}
\begin{array}{rcl}
 \mathsf{patchSource}(Q,\kw{replaceMethodBody}(g,i_f,b))
   &=&Q[\mathit{source}(i_f).\mathit{body}\mapsto b],\\
 \mathsf{patchSource}(Q,\kw{replaceMethodDefinition}(g,M,i_f,f))
   &=&Q[\mathit{source}(i_f)\mapsto f],\\
 \mathsf{patchSource}(Q,\kw{addMethodDefinition}(g,M,f))
   &=&\mathsf{insertMethodSource}(Q,\decl_Q(M),f),\\
 \mathsf{patchSource}(Q,\kw{removeMethodDefinition}(g,M,i_f))
   &=&\mathsf{deleteMethodSource}(Q,\decl_Q(M),i_f),\\
 \mathsf{patchSource}(Q,\kw{replaceSlotDeclarations}(g,M,\delta_s))
   &=&Q[\mathit{source}(\decl_Q(M)).\mathit{slots}\mapsto\delta_s],\\
 \mathsf{patchSource}(Q,\kw{replaceNestedDeclarations}(g,M,\delta_n))
   &=&Q[\mathit{source}(\decl_Q(M)).\mathit{nested}\mapsto\delta_n],\\
 \mathsf{patchSource}(Q,\kw{replaceMixinInitializer}(g,M,\iota))
   &=&Q[\mathit{source}(\decl_Q(M)).\mathit{initializer}\mapsto\iota].
\end{array}
\end{equation}
$\mathsf{insertMethodSource}$ requires a fresh method identity and absent
selector; $\mathsf{deleteMethodSource}$ requires that the identity is a direct
method of the named mixin.  Definition and body replacement preserve the
method identity; delete followed by add cannot reuse it.

Because the transaction is nonconflicting, the fold
\begin{equation}\label{eq:apply-source-commands}
\begin{aligned}
 \mathsf{applySourceCommands}(P,\epsilon)&=P,\\
 \mathsf{applySourceCommands}(P,\delta\cdot\bar\delta)
   &=\mathsf{applySourceCommands}(\mathsf{patchSource}(P,\delta),\bar\delta)
\end{aligned}
\end{equation}
is independent of the chosen canonical locus order.  The patched source
candidate is
\begin{equation}\label{eq:patched-source-candidate}
 Q=\mathsf{applySourceCommands}(P,\mathsf{codeCommands}(\Delta)).
\end{equation}
The ordinary elaborator is rerun on the complete affected dependency closure:
\begin{equation}\label{eq:reflective-reelaboration}
 \begin{aligned}
 &\mathsf{reelaborateProgram}(P,Q)=P'
 \quad\Longleftrightarrow\\[-1mm]
 &\quad P'.\mathit{source}=Q.\mathit{source}\ \land\\
 &\quad P'.\mathit{unit}=Q.\mathit{unit}\ \land\\
 &\quad \forall K\in\dom(P.\mathit{declMixin})\cap\dom(Q.\mathit{source}).\;
    P'.\mathit{declMixin}(K)=P.\mathit{declMixin}(K)\ \land\\
 &\quad \forall f,f'\text{ denoting one surviving declaration}.\;
    \mathsf{methodId}(f')=\mathsf{methodId}(f)\ \land\\
 &\quad \mathsf{coreProgramOK}(P').
 \end{aligned}
\end{equation}
The last clause is Equation~\eqref{eq:annotated-program-validity}; it
recomputes all derived fields of
Equation~\eqref{eq:program-definition}, including synthesized accessors,
lexical parents, scope chains, activation members, owner maps, initializer
normal forms, and method-body maps.  It is undefined on a parse, binding,
access, arity, duplicate-selector, or other static error.  This is why a
transaction cannot expose a half-elaborated program.

\subsection{Class graphs, object layouts, and individual objects}

Let the command projections be
\begin{equation}\label{eq:class-command-projections}
\begin{aligned}
 \mathsf{classCommands}(\Delta)&=
  \mathsf{filter}_{\mathit{ClassCommand}}(\Delta),\\
 \mathsf{objectClassCommands}(\Delta)&=
  \mathsf{filter}_{\mathit{ObjectClassCommand}}(\Delta),\\
 \mathsf{objectSlotCommands}(\Delta)&=
  \mathsf{filter}_{\mathit{ObjectSlotCommand}}(\Delta).
\end{aligned}
\end{equation}
Their heap effects are the distinct partial functions
\begin{equation}\label{eq:class-object-heap-effects}
\begin{aligned}
 \mathsf{applyClassCommand}(H,\kw{changeSuperclass}(g,C,S))
   &=H[C.S\mapsto S],\\
 \mathsf{applyClassCommand}(H,\kw{changeClassEnclosingObject}(g,C,eo))
   &=H[C.eo\mapsto eo],\\
 \mathsf{applyObjectClassCommand}(H,\kw{changeObjectClass}(g,o,C))
   &=H[o.C\mapsto C].
\end{aligned}
\end{equation}
The corresponding phases are
\begin{equation}\label{eq:class-object-command-folds}
\begin{aligned}
 \mathsf{applyClassCommands}(H,\epsilon)&=H,\\
 \mathsf{applyClassCommands}(H,\delta\cdot\bar\delta)
   &=\mathsf{applyClassCommands}
       (\mathsf{applyClassCommand}(H,\delta),\bar\delta),\\
 \mathsf{applyObjectClassCommands}(H,\epsilon)&=H,\\
 \mathsf{applyObjectClassCommands}(H,\delta\cdot\bar\delta)
   &=\mathsf{applyObjectClassCommands}\\
   &\quad(\mathsf{applyObjectClassCommand}(H,\delta),\bar\delta).
\end{aligned}
\end{equation}
The command sequences are in canonical locus order.  Neither phase changes
$H(C).\chi_C$: creation provenance records how the class was created, not its
current superclass or enclosing object.

The three projections in Equation~\eqref{eq:class-command-projections} must
not be conflated.  A class command changes the graph or lexical context shared
by all instances of $C$.  An object-class command changes only $o.C$ and then
subjects that object to layout reconciliation.  An object-slot command leaves
both the class graph and $o.C$ unchanged and writes one slot only after the
final layout has been established.

For every ordinary object, layout reconciliation against the final program
and class graph is
\begin{equation}\label{eq:reconcile-object-layout}
\begin{aligned}
 L_o&=\mathsf{set}(\mathsf{layout}_{P',H_c}(\classof_{H_c}(o))),\\
 \mathsf{reconcileObjectLayout}(P',H_c,o)
   &=H_c[o.\sigma\mapsto
       [j\in L_o\mapsto
          \begin{cases}H_c(o).\sigma(j)&j\in\dom(H_c(o).\sigma),\\
                        \nil&j\notin\dom(H_c(o).\sigma),
          \end{cases}]] .
\end{aligned}
\end{equation}
Thus stable physical slot identities preserve retained state, newly introduced
slots contain $\nil$, and removed slots become unreachable in the same atomic
commit.  The operation applies to every object whose old and new layouts
differ, including every instance affected transitively by a superclass or
mixin-slot edit and an individual object named by \kw{changeObjectClass}.
If $P.\mathit{mixin}(M)=\langle K,\phi,\delta,\nu,\iota\rangle$, define
\begin{equation}\label{eq:nested-declaration-keys}
\begin{aligned}
 \mathsf{ownNestedKeys}_{P}(M)&=\dom(\nu),\\
 \mathsf{nestedKeys}_{P,H}(C)
  &=\bigcup_{R\in\mathsf{set}(\mathsf{chain}_H(C))\setminus\{\Top\}}
       \mathsf{ownNestedKeys}_{P}(\mixin_H(R)).
\end{aligned}
\end{equation}
Nested-class caches are reconciled separately using
Equation~\eqref{eq:nested-declaration-keys}:
\begin{equation}\label{eq:reconcile-nested-cache}
 \mathsf{reconcileNestedCache}(P',H,o)=
 H[o.\eta\mapsto H(o).\eta\mathbin{\upharpoonright}
   \mathsf{nestedKeys}_{P',H}(\classof_H(o))].
\end{equation}
Removing a declaration therefore removes its memoized reachability; adding one
does not eagerly create a class.  Existing class objects remain ordinary live
objects if referenced elsewhere.

Using the VM identity order $\prec$ fixed by
Equation~\eqref{eq:all-instances-of}, define
\begin{equation}\label{eq:reconciliation-phases}
\begin{aligned}
 \mathsf{objectIds}(H)
   &=\mathsf{sort}_{\prec}
       \{o\in\dom(H)\mid H(o)=\kw{object}\langle C,\sigma,\eta\rangle\},\\
 \mathsf{reconcileObjectLayouts}(P',H,\epsilon)&=H,\\
 \mathsf{reconcileObjectLayouts}(P',H,o\cdot\bar O)
   &=\mathsf{reconcileObjectLayouts}\\
   &\quad(P',\mathsf{reconcileObjectLayout}(P',H,o),\bar O),\\
 \mathsf{reconcileAllObjectLayouts}(P',H)
   &=\mathsf{reconcileObjectLayouts}(P',H,\mathsf{objectIds}(H)),\\
 \mathsf{reconcileNestedCaches}(P',H,\epsilon)&=H,\\
 \mathsf{reconcileNestedCaches}(P',H,o\cdot\bar O)
   &=\mathsf{reconcileNestedCaches}\\
   &\quad(P',\mathsf{reconcileNestedCache}(P',H,o),\bar O),\\
 \mathsf{reconcileAllNestedCaches}(P',H)
   &=\mathsf{reconcileNestedCaches}(P',H,\mathsf{objectIds}(H)).
\end{aligned}
\end{equation}

After all reconciliations, an explicit mirror slot write is
\begin{equation}\label{eq:reflective-object-slot-write}
 \begin{aligned}
 &\mathsf{applyObjectSlotWrite}(P',H,\kw{writeObjectSlot}(g,o,j,v))
    =H[o.\sigma(j)\mapsto v]\\
 &\hspace{28mm}\text{if }j\in
    \mathsf{set}(\mathsf{layout}_{P',H}(\classof_H(o))).
 \end{aligned}
\end{equation}
It is undefined otherwise.  This operation is different from the base setter
and from initializer direct write; no source-level access check or synthesized
method is involved.  All actor-isolation checks are nevertheless retained.

The object-slot phase is
\begin{equation}\label{eq:object-slot-command-fold}
\begin{aligned}
 \mathsf{applyObjectSlotWrites}(P',H,\epsilon)&=H,\\
 \mathsf{applyObjectSlotWrites}(P',H,\delta\cdot\bar\delta)
   &=\mathsf{applyObjectSlotWrites}
       (P',\mathsf{applyObjectSlotWrite}(P',H,\delta),\bar\delta).
\end{aligned}
\end{equation}

\subsection{Activations and continuations}

The command projections for the remaining two phases are
\begin{equation}\label{eq:control-command-projections}
\begin{aligned}
 \mathsf{activationCommands}(\Delta)&=
  \mathsf{filter}_{\mathit{ActivationCommand}}(\Delta),\\
 \mathsf{continuationCommands}(\Delta)&=
  \mathsf{filter}_{\mathit{ContinuationCommand}}(\Delta).
\end{aligned}
\end{equation}

For an activation command other than continuation transfer, define
\begin{equation}\label{eq:activation-field-effects}
\begin{aligned}
 \mathsf{applyActivationCommand}(H,\kw{writeActivationParameter}(g,a,x,v))
   &=H[a.\rho(x)\mapsto v],\\
 \mathsf{applyActivationCommand}(H,\kw{writeActivationLocal}(g,a,j,v))
   &=H[a.\ell(j)\mapsto v],\\
 \mathsf{applyActivationCommand}(H,\kw{changeActivationCurrentClass}(g,a,C))
   &=H[a.C\mapsto C],\\
 \mathsf{applyActivationCommand}(H,\kw{changeActivationContinuation}(g,a,k))
   &=H[a.k\mapsto k],\\
 \mathsf{applyActivationCommand}(H,\kw{markActivationUncontinuable}(g,a))
   &=\mathsf{markUncontinuable}(H,a).
\end{aligned}
\end{equation}
The first case is reflectively permitted even though parameter slots are
immutable at base level.  The second requires a physical local identity
already present in $H(a).\ell$.  A changed current class must remain on the
current instance's class chain.  A changed continuation is either $\nil$ or a
strictly lower live stack activation, so continuation links remain acyclic.
Marking uses Equations~\eqref{eq:dependents} and
\eqref{eq:mark-uncontinuable}.

The activation phase is
\begin{equation}\label{eq:activation-command-fold}
\begin{aligned}
 \mathsf{applyActivationCommands}(H,\epsilon)&=H,\\
 \mathsf{applyActivationCommands}(H,\delta\cdot\bar\delta)
   &=\mathsf{applyActivationCommands}
       (\mathsf{applyActivationCommand}(H,\delta),\bar\delta).
\end{aligned}
\end{equation}

Equation~\eqref{eq:resume} supplies the stack operation used by reflective
continuation transfer.  The corresponding state transformer is
\begin{equation}\label{eq:reflective-continuation-transfer}
\begin{aligned}
 &\delta_c=\kw{continueAtActivation}(g,A,a,v),\\
 &\mathsf{applyContinuationTransfer}(H,\mathcal X,\delta_c)
      =\langle H',\mathcal X'\rangle,\\
 &\mathcal X(A)=\kw{runningTurn}\langle\mathcal A,t,p,u\rangle,\qquad
 H(a).q_a=\kw{true},\qquad
 \mathsf{resume}(\mathcal A,a,v)=\langle\mathcal A',v\rangle,\\
 &H'=H[a.k\mapsto\nil],\qquad
 \mathcal X'=\mathcal X[A\mapsto
       \kw{runningTurn}\langle\mathcal A',v,p,u\rangle].
\end{aligned}
\end{equation}
It is undefined for an idle actor, a retained but no longer live activation,
or an uncontinuable activation.  At most one continuation transfer per actor
may occur in a transaction.  The discarded stack suffix remains in the heap,
as in ordinary return, so debuggers may retain it; its continuation dependents
are marked uncontinuable before commit.  This gives reflection continuation
power without introducing a second, incompatible continuation representation.

The continuation-transfer phase is
\begin{equation}\label{eq:continuation-command-fold}
\begin{aligned}
 \mathsf{applyContinuationTransfers}(H,\mathcal X,\epsilon)
   &=\langle H,\mathcal X\rangle,\\
 \mathsf{applyContinuationTransfers}
      (H,\mathcal X,\delta\cdot\bar\delta)
   &=\mathsf{applyContinuationTransfers}(H',\mathcal X',\bar\delta)\\
 &\quad\text{where }\langle H',\mathcal X'\rangle=
       \mathsf{applyContinuationTransfer}(H,\mathcal X,\delta).
\end{aligned}
\end{equation}

\subsection{Debugger stack installation and full-speed resumption}

The debugger-command projection is
\begin{equation}\label{eq:debugger-command-projection}
 \mathsf{debuggerCommands}(\Delta)=
   \mathsf{filter}_{\mathit{DebuggerCommand}}(\Delta).
\end{equation}
For a template $\mathcal S=\langle\psi_1,\ldots,\psi_m\rangle$, key
resolution is the partial function
\begin{equation}\label{eq:resolve-stack-template-keys}
\begin{aligned}
 \mathsf{resolveTemplateKeys}(H,\mathcal S)&=\nu,\\
 \nu(\kw{retain}(a))&=a
   &&\text{only if }H(a)\text{ is an activation record},\\
 \nu(\kw{freshFrame}(u_i))&=a_i
   &&a_i\text{ is fresh},\\
 \nu(\eta_i)&\ne\nu(\eta_j)&&i\ne j.
\end{aligned}
\end{equation}
Fresh choices are simultaneous.  Substitution $\nu(R)$ replaces every
template-key occurrence in a continuation, home, or lexical environment by
its resolved activation identity.

The well-formedness predicate for a resolved stack image is
{\footnotesize
\begin{equation}\label{eq:debug-stack-image-validity}
\begin{split}
 \mathsf{debugStackOK}(P,H,
    \langle a_1,\kappa_1\rangle\cdots\langle a_m,\kappa_m\rangle)
 \quad\Longleftrightarrow\quad{}
 &m>0\ \land\ \mathsf{pairwiseDistinct}(a_1,\ldots,a_m)\ \land\\
 &\mathsf{decodeControlImage}
    (\langle a_1,\kappa_1\rangle\cdots\langle a_m,\kappa_m\rangle)
       \downarrow\ \land\\
 &\forall i\le m.\ H(a_i)\text{ is an activation record}\ \land
       H(a_i).q_a=\kw{true}\ \land\\
 &\forall i\le m.\ \mathsf{controlLocation}(\kappa_i)\downarrow
       \ \land\ \mathsf{controlRefs}(\kappa_i)\subseteq\dom(H)\ \land\\
 &\forall i\le m.\ H(a_i).k=\nil\ \lor\
       \exists j<i.\ H(a_i).k=a_j\ \land\\
 &\forall i\le m.\ H(a_i).h=\nil\ \lor\ H(a_i).h\in\dom(H).
\end{split}
\end{equation}
}
$\mathsf{controlRefs}$ is the instance of the structural reference extractor
of Equation~\eqref{eq:structural-reference-extractor} for terms and
intra-activation frames.  The control-map definedness premise validates a
source or machine location; a bytecode platform additionally requires the operand-stack and
temporary layout encoded by $\pi_i$ to match the stack map at the corresponding
PC.  This is the implementation representation of the same predicate, not a
weaker alternate rule.

Removed live activations remain available to retained mirrors but cease to be
continuable:
\begin{equation}\label{eq:retire-removed-stack-activations}
 \mathsf{retireRemoved}(H,R)=
 H[\,a.k\mapsto\nil,\ a.q_a\mapsto\kw{false}\mid a\in R\,].
\end{equation}
Define atomic template materialization by
\begin{equation}\label{eq:materialize-stack-template}
\begin{split}
 &\mathsf{materializeStackTemplate}(P,H,\mathcal A,\mathcal S)
      =\langle H',\mathcal A',t'\rangle
 \quad\Longleftrightarrow\\
 &\quad\nu=\mathsf{resolveTemplateKeys}(H,\mathcal S)\ \land\\
 &\quad\mathcal K=
   \langle\nu(\psi_1.\eta),\psi_1.\kappa\rangle\cdots
   \langle\nu(\psi_m.\eta),\psi_m.\kappa\rangle\ \land\\
 &\quad\mathsf{decodeControlImage}(\mathcal K)
      =\langle\mathcal A',t'\rangle\ \land\\
 &\quad R=\mathsf{set}(\mathsf{activationIds}(\mathcal A))
       \setminus\mathsf{set}(\mathsf{activationIds}(\mathcal A'))\ \land\\
 &\quad H_0=\mathsf{retireRemoved}(H,R)\ \land\\
 &\quad H'=H_0[\,\nu(\psi_i.\eta)\mapsto\nu(\psi_i.R)
          \mid1\le i\le m\,]\ \land\\
 &\quad\mathsf{debugStackOK}(P,H',\mathcal K).
\end{split}
\end{equation}
It is undefined if any key, record, control point, reference, continuation, or
stack-shape condition fails.  Since only $H'$ and the reconstructed
$\langle\mathcal A',t'\rangle$ leave the function, neither allocation nor
retirement is visible on failure.

The effect of one debugger command, issued by actor $A$, is the following
partial function.  Pausing allocates a fresh token:
\begin{equation}\label{eq:pause-actor-effect}
\begin{aligned}
 &\mathsf{applyDebuggerCommand}
   (P,H,\mathcal X,A,\kw{pauseActor}(g,B,d_s))
       =\langle H,\mathcal X'\rangle,\\
 &\mathcal X(B)=\kw{runningTurn}\langle\mathcal A,t,p,u\rangle,
 \quad A\ne B,\quad\mathcal A\ne\epsilon,\quad r_s\text{ fresh},\\
 &\mathcal X'=\mathcal X[B\mapsto
   \kw{pausedTurn}\langle\mathcal A,t,p,u,r_s,d_s\rangle].
\end{aligned}
\end{equation}
A debugger running in the actor that it is editing uses the nonreturning
current-stack replacement:
\begin{equation}\label{eq:replace-current-actor-stack}
\begin{aligned}
 &\mathsf{applyDebuggerCommand}
   (P,H,\mathcal X,A,\kw{replaceCurrentActorStack}(g,B,\mathcal S))
       =\langle H',\mathcal X'\rangle,\\
 &A=B,\quad
 \mathcal X(B)=\kw{runningTurn}\langle\mathcal A,t,p,u\rangle,\\
 &\mathsf{materializeStackTemplate}(P,H,\mathcal A,\mathcal S)
       =\langle H',\mathcal A',t'\rangle,\\
 &\mathcal X'=\mathcal X[B\mapsto
       \kw{runningTurn}\langle\mathcal A',t',p,u\rangle].
\end{aligned}
\end{equation}
The reflective invocation that requested this command is among the discarded
or replaced frames, so it does not return.  This is the operation used by an
in-actor debugger after a top-level resumable-exception handler has constructed
the desired continuation.
For the common resumable-exception case, the derived constructor is
\begin{equation}\label{eq:resume-exception-stack-template}
 \mathsf{resumeExceptionStack}
  (\mathcal S_0\cdot\langle\psi_f\rangle\cdot\mathcal S_d,a_f,v)
 =\mathcal S_0\cdot
   \langle\mathsf{promoteFrame}(\psi_f,v)\rangle
 \quad\text{if }\psi_f.\eta=\kw{retain}(a_f).
\end{equation}
It removes the handler and debugger suffix $\mathcal S_d$ and makes $v$ the
result of the original \kws{signal} send.  Submitting this template through
Equation~\eqref{eq:replace-current-actor-stack} therefore resumes precisely at
the protected failure point under the normal VM relation.

For a paused actor, replacement checks the exact pause token and issues a new
one:
\begin{equation}\label{eq:replace-paused-actor-stack}
\begin{aligned}
 &\mathsf{applyDebuggerCommand}
  (P,H,\mathcal X,A,\kw{replacePausedActorStack}(g,B,r_s,\mathcal S))
      =\langle H',\mathcal X'\rangle,\\
 &\mathcal X(B)=\kw{pausedTurn}\langle
      \mathcal A,t,p,u,r_s,d_s\rangle,\\
 &\mathsf{materializeStackTemplate}(P,H,\mathcal A,\mathcal S)
      =\langle H',\mathcal A',t'\rangle,\quad r_s'\text{ fresh},\\
 &\mathcal X'=\mathcal X[B\mapsto\kw{pausedTurn}\langle
      \mathcal A',t',p,u,r_s',d_s\rangle].
\end{aligned}
\end{equation}
Resuming an unchanged paused stack selects ordinary VM execution:
\begin{equation}\label{eq:resume-paused-actor-full-speed}
\begin{aligned}
 &\mathsf{applyDebuggerCommand}
  (P,H,\mathcal X,A,\kw{resumePausedActorAtFullSpeed}(g,B,r_s))
      =\langle H,\mathcal X'\rangle,\\
 &\mathcal X(B)=\kw{pausedTurn}\langle
      \mathcal A,t,p,u,r_s,d_s\rangle,\\
 &\mathcal X'=\mathcal X[B\mapsto
      \kw{runningTurn}\langle\mathcal A,t,p,u\rangle].
\end{aligned}
\end{equation}
Replacement and full-speed resumption may instead be one atomic operation:
\begin{equation}\label{eq:replace-resume-paused-actor-stack}
\begin{aligned}
 &\mathsf{applyDebuggerCommand}(P,H,\mathcal X,A,
   \kw{replaceAndResumePausedActorStack}(g,B,r_s,\mathcal S))
      =\langle H',\mathcal X'\rangle,\\
 &\mathcal X(B)=\kw{pausedTurn}\langle
      \mathcal A,t,p,u,r_s,d_s\rangle,\\
 &\mathsf{materializeStackTemplate}(P,H,\mathcal A,\mathcal S)
      =\langle H',\mathcal A',t'\rangle,\\
 &\mathcal X'=\mathcal X[B\mapsto
      \kw{runningTurn}\langle\mathcal A',t',p,u\rangle].
\end{aligned}
\end{equation}
``Full speed'' is not a second language semantics: after either resume equation,
Rule \rulelabel{Actor-Sequential} again lifts the ordinary sequential relation.
An interpreter, baseline VM, or optimizing VM conforms by implementing that
same relation from the installed control image.

The debugger phase is the explicit fold
\begin{equation}\label{eq:debugger-command-fold}
\begin{aligned}
 \mathsf{applyDebuggerCommands}(P,H,\mathcal X,A,\epsilon)
   &=\langle H,\mathcal X\rangle,\\
 \mathsf{applyDebuggerCommands}
   (P,H,\mathcal X,A,\delta\cdot\bar\delta)
   &=\mathsf{applyDebuggerCommands}
       (P,H',\mathcal X',A,\bar\delta)\\
 &\quad\text{where }\langle H',\mathcal X'\rangle={}\\[-1mm]
 &\hspace{29mm}\mathsf{applyDebuggerCommand}(P,H,\mathcal X,A,\delta).
\end{aligned}
\end{equation}

\paragraph{Simulator correspondence.}
Let $D$ be a Simulator state represented by ordinary Newspeak objects and let
$D\approx_{\mathsf{sim}}\Sigma$ mean that its bytecode frames and program
counters decode to the stack and control image of
Equation~\eqref{eq:stack-control-image}.  Write
$\longmapsto_{\mathsf{Simulator}}^{+}$ for a nonempty finite Simulator
execution and $\longmapsto^{*}$ for a possibly empty finite semantic
execution.  An admitted compiler/Simulator pair supplies the three-part
certificate
\begin{equation}\label{eq:simulator-step-correspondence}
\begin{aligned}
 D\approx_{\mathsf{sim}}\Sigma\ \land\
 D\longmapsto_{\mathsf{Simulator}}D'
 &\Rightarrow
 \exists\Sigma'.\ \Sigma\longmapsto^{*}\Sigma'
       \land D'\approx_{\mathsf{sim}}\Sigma',\\
 D\approx_{\mathsf{sim}}\Sigma\ \land\ \Sigma\longmapsto\Sigma'
 &\Rightarrow
 \exists D'.\ D\longmapsto_{\mathsf{Simulator}}^{+}D'
       \land D'\approx_{\mathsf{sim}}\Sigma',\\
 D\approx_{\mathsf{sim}}\Sigma
 &\Rightarrow \mathsf{resumeAtFullSpeed}(D)=\Sigma.
\end{aligned}
\end{equation}
The first clause permits bytecode-level stuttering while proving instruction
soundness; the second prevents the Simulator from omitting a semantic step;
the third closes the critical boundary back to native execution without a
change of represented configuration.  The certificate lifts the first clause
from one instruction to every finite Simulator execution.
The Simulator can inspect, push, pop, and copy its ordinary frame-image values;
Equations~\eqref{eq:replace-current-actor-stack} and
\eqref{eq:replace-resume-paused-actor-stack} cross back to normal VM execution.
Proving Equation~\eqref{eq:simulator-step-correspondence} for the actual
Simulator and bytecode set is an implementation-refinement proof; the language
semantics fixes the state and transfer contract it must implement.

\subsection{Preparation, validation, and atomic commit}

All phase ordering is semantic, not user-visible.  Define the partial
preparation function
\begin{equation}\label{eq:prepare-reflection}
\begin{split}
 &\mathsf{prepareReflection}(\Omega,A,\Delta)=\Omega'\\[-1mm]
 &\quad\Longleftrightarrow\quad
 P'=\mathsf{reelaborateProgram}(P,Q)\ \land\\
 &H_1=\mathsf{applyClassCommands}
       (\mathsf{vmStore}(\Omega),\mathsf{classCommands}(\Delta))\ \land\\
 &H_c=\mathsf{applyObjectClassCommands}
       (H_1,\mathsf{objectClassCommands}(\Delta))\ \land\\
 &H_l=\mathsf{reconcileAllObjectLayouts}(P',H_c)\ \land\\
 &H_n=\mathsf{reconcileAllNestedCaches}(P',H_l)\ \land\\
 &H_o=\mathsf{applyObjectSlotWrites}
       (P',H_n,\mathsf{objectSlotCommands}(\Delta))\ \land\\
 &H_a=\mathsf{applyActivationCommands}
       (H_o,\mathsf{activationCommands}(\Delta))\ \land\\
 &\langle H_d,\mathcal X_d\rangle=
     \mathsf{applyDebuggerCommands}
       (P',H_a,\mathcal X,A,\mathsf{debuggerCommands}(\Delta))\ \land\\
 &\langle H_t,\mathcal X'\rangle=
     \mathsf{applyContinuationTransfers}
       (H_d,\mathcal X_d,\mathsf{continuationCommands}(\Delta))\ \land\\
 &\Omega_0=\Omega[P\mapsto P',\mathcal X\mapsto\mathcal X',n\mapsto n+1]\ \land\\
 &\Omega'=\mathsf{repartitionVMStore}(\Omega_0,H_t).
\end{split}
\end{equation}
The fresh activations introduced by stack installation are
\begin{equation}\label{eq:fresh-installed-activations}
 \mathsf{freshInstalled}_{\Omega}(\mathcal X',B)=
 \mathsf{set}(\mathsf{stackActivationIds}_{\mathcal X'}(B))
 \setminus\dom(\mathsf{vmStore}(\Omega)).
\end{equation}
$\mathsf{repartitionVMStore}$ returns every pre-existing identity to the same
$V$ or $\mathcal H(A)$ domain from which it came, and places each member of
Equation~\eqref{eq:fresh-installed-activations} in $\mathcal H(B)$.  The sets
for distinct actors must be disjoint.  Actor-isolation validation below rejects
any installed record containing an illegal near reference; reflection creates
no cross-actor near reference.

The validation predicates below are over tagged record sets, so their
quantifiers have explicit domains:
\begin{equation}\label{eq:reflective-record-id-sets}
\begin{aligned}
 \mathsf{classRecordIds}(H)
  &=\{C\in\dom(H)\mid
       H(C)=\kw{class}\langle M,S,eo,C_m,q,\chi_C\rangle\},\\
 \mathsf{ordinaryObjectIds}(H)
  &=\{o\in\dom(H)\mid H(o)=\kw{object}\langle C,\sigma,\eta\rangle\},\\
 \mathsf{activationRecordIds}(H)
  &=\{a\in\dom(H)\mid
       H(a)=\kw{activation}\langle p,\rho,\ell,o,C,k,h,q_a,\gamma\rangle\}.
\end{aligned}
\end{equation}
Class-chain termination is the predicate
\begin{equation}\label{eq:reflective-class-chain-validity}
\begin{split}
 \mathsf{reachesTop}_{H}(C)\quad\Longleftrightarrow\quad{}
 &\exists k\in\mathbb N,C_0,\ldots,C_k.\;
 C_0=C\ \land\ C_k=\Top\ \land\\[-1mm]
 &\mathsf{pairwiseDistinct}(C_0,\ldots,C_k)\ \land\\
 &\forall i<k.\;C_i\in\mathsf{classRecordIds}(H)\ \land
                  \super_H(C_i)=C_{i+1}.
\end{split}
\end{equation}
Individual class-record consistency is
\begin{equation}\label{eq:reflective-class-record-validity}
\begin{split}
 \mathsf{classRecordOK}(P,H,C)\quad\Longleftrightarrow\quad{}
 &\mixin_H(C)\in\dom(P.\mathit{mixin})\ \land\\
 &H(C).C_m\in\mathsf{classRecordIds}(H)\ \land\\
 &H(C).eo\in\dom(H)\cup\{\nil\}\ \land
   (C=\Object\ \lor\ \super_H(C)\ne\Top).
\end{split}
\end{equation}
The complete class-graph check is
\begin{equation}\label{eq:reflective-class-graph-validity}
\begin{split}
 \mathsf{classGraphOK}(P,H)\quad\Longleftrightarrow\quad{}
 &\Object\in\mathsf{classRecordIds}(H)\ \land
   \super_H(\Object)=\Top\ \land\\
 &\forall C\in\mathsf{classRecordIds}(H).\;
      \mathsf{reachesTop}_{H}(C)\ \land\\
 &\forall C\in\mathsf{classRecordIds}(H)\setminus\{\Top\}.\;
      \mathsf{classRecordOK}(P,H,C).
\end{split}
\end{equation}
The declaration retained by class provenance and the complete ordinary-object
check are defined by
\begin{equation}\label{eq:class-origin-declaration}
 \mathsf{originDecl}(\kw{expression}(K))=K,
 \qquad \mathsf{originDecl}(\kw{actor}(K))=K.
\end{equation}
The cached class for one nested declaration is valid exactly when
\begin{equation}\label{eq:cached-nested-class-validity}
\begin{split}
 \mathsf{cachedNestedClassOK}(H,o,d)\quad\Longleftrightarrow\quad{}
 &H(o).\eta(d)\in\mathsf{classRecordIds}(H)\ \land\\
 &H(H(o).\eta(d)).eo=o\ \land\\
 &\mathsf{originDecl}(H(H(o).\eta(d)).\chi_C)=d.
\end{split}
\end{equation}
An ordinary object record is valid when
\begin{equation}\label{eq:reflective-object-record-validity}
\begin{aligned}
 \mathsf{objectRecordOK}(P,H,o)\quad\Longleftrightarrow\quad{}
 &\classof_H(o)\in\mathsf{classRecordIds}(H)\ \land\\
 &\dom(H(o).\sigma)=
   \mathsf{set}(\mathsf{layout}_{P,H}(\classof_H(o)))\ \land\\
 &\dom(H(o).\eta)\subseteq
   \mathsf{nestedKeys}_{P,H}(\classof_H(o))\ \land\\
 &\forall d\in\dom(H(o).\eta).\;\mathsf{cachedNestedClassOK}(H,o,d).
\end{aligned}
\end{equation}
The store-wide object check is
\begin{equation}\label{eq:reflective-object-validity}
 \mathsf{objectStoreOK}(P,H)\quad\Longleftrightarrow\quad
 \forall o\in\mathsf{ordinaryObjectIds}(H).\;
   \mathsf{objectRecordOK}(P,H,o).
\end{equation}

For activation validation, first extract the activation identities from each
actor's live stack:
\begin{equation}\label{eq:live-activation-order}
\begin{gathered}
 \mathsf{activationIds}(\epsilon)=\epsilon,\\
 \mathsf{activationIds}(\mathcal A\mathbin{\cdot}\langle a,\pi\rangle)
   =\mathsf{activationIds}(\mathcal A)\mathbin{\cdot}\langle a\rangle,\\
 \mathsf{stackActivationIds}_{\mathcal X}(A)=\epsilon
   \quad\text{if }\mathcal X(A)=\kw{idle},\\
 \mathsf{stackActivationIds}_{\mathcal X}(A)=\mathsf{activationIds}(\mathcal A),\\
 \hfill\text{if }\mathcal X(A)=\kw{runningTurn}\langle\mathcal A,t,p,u\rangle,\\
 \mathsf{stackActivationIds}_{\mathcal X}(A)=\mathsf{activationIds}(\mathcal A),\\
 \hfill\text{if }\mathcal X(A)=
   \kw{pausedTurn}\langle\mathcal A,t,p,u,r_s,d_s\rangle,\\
 \mathsf{liveActivationIds}(\mathcal X)
   =\bigcup_{A\in\dom(\mathcal X)}
       \mathsf{set}(\mathsf{stackActivationIds}_{\mathcal X}(A)),\\
 \mathsf{below}_{\mathcal X}(A,k,a)\Longleftrightarrow\exists i<j.\\
 \mathsf{stackActivationIds}_{\mathcal X}(A)[i]=k\ \land
 \mathsf{stackActivationIds}_{\mathcal X}(A)[j]=a.
\end{gathered}
\end{equation}
Definedness of every partial expression below uses the notation
\begin{equation}\label{eq:definedness-notation}
 e_p\downarrow\ \Longleftrightarrow\ \exists v.\;e_p=v,
 \qquad e_p\uparrow\ \Longleftrightarrow\ \neg(e_p\downarrow).
\end{equation}
Surviving method declarations determine the parameter domain of a live method
activation:
\begin{equation}\label{eq:method-formals}
 \mathsf{methodFormals}_{P}(i_f)=\bar x
 \Longleftrightarrow
 \exists M,s,\alpha,\ell,b,K.\;
 \methods_P(M)(s)=\langle i_f,s,\alpha,\bar x,\ell,b,K\rangle.
\end{equation}
The validity of one live activation is
\begin{equation}\label{eq:live-activation-record-validity}
\begin{aligned}
 \mathsf{liveActivationRecordOK}(P,H,\mathcal X,a)
 &\Longleftrightarrow
 \exists p,\rho,\ell,o,C,k,h,q_a,\gamma.\;\\[-1mm]
 &\quad H(a)=\kw{activation}\langle
       p,\rho,\ell,o,C,k,h,q_a,\gamma\rangle\ \land\\
 &\quad\bigl((p=\kw{top}\ \land\ \rho=\ell=\varnothing\ \land
              o=C=h=\nil)\ \lor\\
 &\qquad(p\ne\kw{top}\ \land\ o\in\dom(H)\ \land
          C\preceq_H\classof_H(o)\ \land\\
 &\hspace{39mm}h\in\mathsf{activationRecordIds}(H))\bigr)\ \land\\
 &\quad\bigl(k=\nil\ \lor
       \exists A.\;\mathsf{below}_{\mathcal X}(A,k,a)\bigr)\ \land\\
 &\quad\bigl(\mathsf{methodFormals}_{P}(p)\uparrow\ \lor\\
 &\hspace{22mm}\dom(\rho)=
       \mathsf{set}(\mathsf{methodFormals}_{P}(p))\bigr).
\end{aligned}
\end{equation}
The aggregate check is
\begin{equation}\label{eq:live-activation-validity}
\begin{split}
 \mathsf{liveActivationsOK}(P,H,\mathcal X)
 \quad\Longleftrightarrow\quad{}
 &\forall a\in\mathsf{liveActivationIds}(\mathcal X).\\[-1mm]
 &\mathsf{liveActivationRecordOK}(P,H,\mathcal X,a).
\end{split}
\end{equation}
Equation~\eqref{eq:live-activation-validity} therefore permits an already
materialized activation of a removed method, but requires the bound parameter
domain to agree with every surviving edited signature.

Let $H_0=\mathsf{vmStore}(\Omega)$ and let $H'=\mathsf{vmStore}(\Omega')$.
The command-specific activation conditions are the following total predicate:
\begin{equation}\label{eq:activation-parameter-edit-validity}
\begin{split}
 &\mathsf{activationEditOK}_{H_0,H',\mathcal X'}
   (\kw{writeActivationParameter}(g,a,x,v))\\
 &\qquad\Longleftrightarrow
   a\in\mathsf{activationRecordIds}(H_0)\ \land x\in\dom(H_0(a).\rho).
\end{split}
\end{equation}
\begin{equation}\label{eq:activation-local-edit-validity}
\begin{split}
 &\mathsf{activationEditOK}_{H_0,H',\mathcal X'}
   (\kw{writeActivationLocal}(g,a,j,v))\\
 &\qquad\Longleftrightarrow
   a\in\mathsf{activationRecordIds}(H_0)\ \land j\in\dom(H_0(a).\ell).
\end{split}
\end{equation}
\begin{equation}\label{eq:activation-class-edit-validity}
\begin{split}
 &\mathsf{activationEditOK}_{H_0,H',\mathcal X'}
   (\kw{changeActivationCurrentClass}(g,a,C))\\
 &\qquad\Longleftrightarrow
   a\in\mathsf{activationRecordIds}(H')\ \land
   C\preceq_{H'}\classof_{H'}(H'(a).o).
\end{split}
\end{equation}
\begin{equation}\label{eq:activation-continuation-edit-validity}
\begin{split}
 &\mathsf{activationEditOK}_{H_0,H',\mathcal X'}
   (\kw{changeActivationContinuation}(g,a,k))\\
 &\qquad\Longleftrightarrow
   a\in\mathsf{liveActivationIds}(\mathcal X')\ \land\\
 &\hspace{34mm}(k=\nil\ \lor\
      \exists A.\;\mathsf{below}_{\mathcal X'}(A,k,a)).
\end{split}
\end{equation}
\begin{equation}\label{eq:activation-uncontinuable-edit-validity}
\begin{split}
 &\mathsf{activationEditOK}_{H_0,H',\mathcal X'}
   (\kw{markActivationUncontinuable}(g,a))\\
 &\qquad\Longleftrightarrow a\in\mathsf{activationRecordIds}(H_0),\\
 &\mathsf{activationEditOK}_{H_0,H',\mathcal X'}(\delta)
   \Longleftrightarrow\kw{true}
   \quad\text{if }\delta\notin\mathit{ActivationCommand}.
\end{split}
\end{equation}
Continuation transfer is already partial in
Equation~\eqref{eq:reflective-continuation-transfer}; its transaction-level
cardinality and the remaining activation edits are checked by
\begin{equation}\label{eq:continuation-transfer-count}
 \mathsf{transferCount}(\Delta,B)=
 \left|\{\delta\in\Delta\mid\exists g,a,v.\;
       \delta=\kw{continueAtActivation}(g,B,a,v)\}\right|.
\end{equation}
\begin{equation}\label{eq:activation-transaction-validity}
\begin{aligned}
 \mathsf{activationEditsOK}(\Omega,\Delta,\Omega')
 \quad\Longleftrightarrow\quad{}
 &\bigl(\forall\delta\in\Delta.\;
   \mathsf{activationEditOK}_{H_0,H',\mathcal X'}(\delta)\bigr)\ \land\\
 &\forall B\in\dom(\mathcal X').\;
   \mathsf{transferCount}(\Delta,B)\le1.
\end{aligned}
\end{equation}

Let $\mathsf{refs}$ be the structural reference extractor, defined for every
runtime constructor, finite map, and sequence by
\begin{equation}\label{eq:structural-reference-extractor}
\begin{aligned}
 \mathsf{refs}(x)&=\{x\}&&x\in\mathit{Ref},\\
 \mathsf{refs}(z)&=\varnothing&&z\text{ is atomic and }z\notin\mathit{Ref},\\
 \mathsf{refs}(\mathsf{ctor}(z_1,\ldots,z_k))
   &=\bigcup_{i=1}^{k}\mathsf{refs}(z_i),\\
 \mathsf{refs}([z_i\mapsto z_i']_{i\in I})
   &=\bigcup_{i\in I}(\mathsf{refs}(z_i)\cup\mathsf{refs}(z_i')).
\end{aligned}
\end{equation}
The actor-residency check is then
\begin{equation}\label{eq:reflective-actor-isolation}
\begin{aligned}
 \mathsf{resident}_{\Omega}(A)
  &=\dom(V)\cup\dom(\mathcal H(A))\cup\dom(\Pi)\cup\dom(\Phi),\\
 \mathsf{actorIsolationOK}(\Omega)\quad\Longleftrightarrow\quad{}
 &\forall A\in\dom(\mathcal H).\;
   \mathsf{refs}(\mathcal X(A))\subseteq\mathsf{resident}_{\Omega}(A)\ \land\\
 &\quad\forall x\in\dom(\mathcal H(A)).\;
   \mathsf{refs}(\mathcal H(A)(x))\subseteq\mathsf{resident}_{\Omega}(A)\ \land\\
 &\forall x\in\dom(V).\;
   \mathsf{refs}(V(x))\subseteq\dom(V)\ \land
   \mathsf{value}_{P,\mathsf{vmStore}(\Omega)}(x).
\end{aligned}
\end{equation}

The finite validation-kind domain, its interpretation, and its fixed order are
{\small
\begin{equation}\label{eq:reflection-validation-kinds}
\begin{gathered}
 \kappa_r::=\kw{classGraphCheck}\mid\kw{objectStoreCheck}
  \mid\kw{liveActivationCheck}\mid\kw{activationEditCheck}
  \mid\kw{actorIsolationCheck},\\
 \kw{classGraphCheck}\mathrel{\triangleleft_{\mathsf{check}}}
 \kw{objectStoreCheck}\mathrel{\triangleleft_{\mathsf{check}}}
 \kw{liveActivationCheck},\\
 \kw{liveActivationCheck}\mathrel{\triangleleft_{\mathsf{check}}}
 \kw{activationEditCheck}\mathrel{\triangleleft_{\mathsf{check}}}
 \kw{actorIsolationCheck},\\
 \mathsf{validationHolds}(\kw{classGraphCheck},\Omega,\Delta,\Omega')
  \Longleftrightarrow\mathsf{classGraphOK}(P',H'),\\
 \mathsf{validationHolds}(\kw{objectStoreCheck},\Omega,\Delta,\Omega')
  \Longleftrightarrow\mathsf{objectStoreOK}(P',H'),\\
 \mathsf{validationHolds}(\kw{liveActivationCheck},\Omega,\Delta,\Omega')
  \Longleftrightarrow\mathsf{liveActivationsOK}(P',H',\mathcal X'),\\
 \mathsf{validationHolds}(\kw{activationEditCheck},\Omega,\Delta,\Omega')
  \Longleftrightarrow\mathsf{activationEditsOK}(\Omega,\Delta,\Omega'),\\
 \mathsf{validationHolds}(\kw{actorIsolationCheck},\Omega,\Delta,\Omega')
  \Longleftrightarrow\mathsf{actorIsolationOK}(\Omega').
\end{gathered}
\end{equation}
}
Define the failed validation set by
\begin{equation}\label{eq:failed-reflection-validations}
 \mathsf{failedValidations}(\Omega,\Delta,\Omega')=
 \{\kappa_r\mid
   \neg\mathsf{validationHolds}(\kappa_r,\Omega,\Delta,\Omega')\}.
\end{equation}
For $\Delta=\langle\delta_1,\ldots,\delta_m\rangle$, also define
\begin{equation}\label{eq:reflection-error-index-sets}
\begin{aligned}
 \mathsf{unauthorizedIndices}(\Omega,A,\Delta)
  &=\{i\mid1\le i\le m\ \land
       \neg\mathsf{authorized}_{\Omega,A}(\delta_i)\},\\
 \mathsf{conflictPairs}(\Delta)
  &=\{(i,j)\mid\mathsf{conflict}_{\Delta}(i,j)\}.
\end{aligned}
\end{equation}
The validity predicate now has no informal conjuncts:
\begin{equation}\label{eq:reflection-validity}
\begin{aligned}
 \mathsf{validReflection}(\Omega,A,\Delta,\Omega')
 \quad\Longleftrightarrow\quad{}
 &\mathsf{prepareReflection}(\Omega,A,\Delta)=\Omega'\ \land\\
 &\mathsf{unauthorizedIndices}(\Omega,A,\Delta)=\varnothing\ \land\\
 &\mathsf{conflictPairs}(\Delta)=\varnothing\ \land\\
 &\mathsf{failedValidations}(\Omega,\Delta,\Omega')=\varnothing.
\end{aligned}
\end{equation}

Successful preparation commits in one VM step:
\[
\dfrac{\mathsf{validReflection}(\Omega,A,\Delta,\Omega')}
 {\Omega\xrightarrow{\;\mathsf{commitReflection}(A,\Delta,n+1)\;}\Omega'}
\quad(\rulelabel{Reflection-Commit})
\]
The rejection diagnostic domain is
\begin{equation}\label{eq:reflection-failure-domain}
\begin{aligned}
 \xi_r::={}&\kw{unauthorizedCommand}(i)
 \mid\kw{commandConflict}(i,j)
 \mid\kw{sourcePatchFailure}\mid\kw{elaborationFailure}\\
 &\mid\kw{statePreparationFailure}\mid\kw{invariantFailure}(\kappa_r).
\end{aligned}
\end{equation}
For compactness, define the two partial candidates
\begin{equation}\label{eq:reflection-diagnostic-candidates}
\begin{aligned}
 \mathsf{reflectionSourceCandidate}(P,\Delta)
  &=\mathsf{applySourceCommands}(P,\mathsf{codeCommands}(\Delta)),\\
 \mathsf{reflectionProgramCandidate}(P,\Delta)
  &=\mathsf{reelaborateProgram}
      (P,\mathsf{reflectionSourceCandidate}(P,\Delta)).
\end{aligned}
\end{equation}
Writing
$\mathcal U_r=\mathsf{unauthorizedIndices}(\Omega,A,\Delta)$ and
$\mathcal C_r=\mathsf{conflictPairs}(\Delta)$, the previously implicit failure
function is the following partial function:
{\small
\begin{equation}\label{eq:reflection-failure}
\mathsf{reflectionFailure}(\Omega,A,\Delta)=
\begin{cases}
 \kw{unauthorizedCommand}(\min\mathcal U_r)
   &\mathcal U_r\ne\varnothing,\\
 \kw{commandConflict}(i,j)
   &\begin{gathered}\mathcal U_r=\varnothing\ \land\
       \mathcal C_r\ne\varnothing\ \land\\
       (i,j)=\min_{\rm lex}\mathcal C_r,
    \end{gathered}\\
 \kw{sourcePatchFailure}
   &\begin{gathered}\mathcal U_r=\mathcal C_r=\varnothing\ \land\\
    \mathsf{reflectionSourceCandidate}(P,\Delta)\uparrow,
    \end{gathered}\\
 \kw{elaborationFailure}
   &\begin{gathered}\mathcal U_r=\mathcal C_r=\varnothing\ \land\\
    \mathsf{reflectionSourceCandidate}(P,\Delta)\downarrow\ \land\\
    \mathsf{reflectionProgramCandidate}(P,\Delta)\uparrow,
    \end{gathered}\\
 \kw{statePreparationFailure}
   &\begin{gathered}\mathcal U_r=\mathcal C_r=\varnothing\ \land\\
    \mathsf{reflectionProgramCandidate}(P,\Delta)\downarrow\ \land\\
    \mathsf{prepareReflection}(\Omega,A,\Delta)\uparrow,
    \end{gathered}\\
 \kw{invariantFailure}(\min_{\triangleleft_{\mathsf{check}}}\mathcal B_r)
   &\begin{gathered}\mathcal U_r=\mathcal C_r=\varnothing\ \land\\
    \mathsf{prepareReflection}(\Omega,A,\Delta)=\Omega'\ \land\\
    \mathcal B_r=\mathsf{failedValidations}(\Omega,\Delta,\Omega')
       \ne\varnothing.
    \end{gathered}
\end{cases}
\end{equation}
}
It is undefined exactly when
$\mathsf{validReflection}(\Omega,A,\Delta,\Omega')$ holds for some $\Omega'$.
Failure leaves the state identical:
\[
\dfrac{\mathsf{reflectionFailure}(\Omega,A,\Delta)=\xi}
 {\Omega\xrightarrow{\;\mathsf{rejectReflection}(A,\Delta,\xi)\;}\Omega}
\quad(\rulelabel{Reflection-Reject})
\]
The case order in Equation~\eqref{eq:reflection-failure} affects only the
diagnostic; it cannot expose an intermediate state.  Except for a successful
Equation~\eqref{eq:replace-current-actor-stack}, whose requesting frame has
itself been replaced, the mirror library reifies the commit receipt or failure
as the return or exception of the primitive mirror message.  Current-stack
replacement instead continues directly from the installed top control.

\subsection{Causal connection, turns, and other VMs}

A reflective transition is a VM safepoint transition: it may occur between any
two actor, delivery, or sequential steps, but never inside one such step.  It
replaces the single $P$ in Equation~\eqref{eq:actor-configuration}; hence every
actor in the cohort observes the same version $n+1$ immediately.  The following
definition fixes the boundary between old and new behavior:
\begin{equation}\label{eq:reflection-version-observation}
\begin{array}{c|c}
\text{computation state at commit}&\text{program observed afterward}\\ \hline
\text{already materialized control }t\text{ and frames }\pi
  &\text{their existing terms, code image, and control map}\\
\text{future method resolution or invocation}&P'\\
\begin{gathered}\text{mailbox or network packet}\\[-1mm]
 \text{not yet dispatched}\end{gathered}&P'\text{ at its later turn}\\
\text{closure object already allocated}&\text{its captured body in }H(c)\\
\text{class or object retaining stable }M,C,i_f&
  \text{the updated records addressed by those identities.}
\end{array}
\end{equation}
Thus an install never splices a new expression into a currently executing
small step.  A method invocation that has already entered its body continues
that body; the next dispatch through the same stable method identity uses its
new definition.  In-flight asynchronous messages carry selectors and
arguments, not resolved methods, so they naturally dispatch against the
version current when Rule \rulelabel{Turn-Start} begins their turn.

Let a distributed system be
\begin{equation}\label{eq:vm-family}
 \mathfrak F:\mathit{VMId}\rightharpoonup\mathit{ActorConfiguration},
 \qquad \mathfrak F(w).w=w.
\end{equation}
Rule \rulelabel{Reflection-Commit} changes exactly one component
$\mathfrak F(w)$.  No rule copies $P'$ into $\mathfrak F(w')$ for $w'\ne w$.
Values in $V$ are shared by every actor of $w$, so their behavior changes with
the cohort's code; a value transferred to another VM is copied under
Equation~\eqref{eq:value-transfer-contract} and thereafter uses the recipient
VM's independently versioned program.  A remote VM may be changed only by a
separate, authorized mirror request delivered to that VM.  Its delivery is
ordered like any other eventual message; the resulting local commit is atomic
at the recipient safepoint.

The same rule governs remote debugging.  An actor mirror transferred to
another actor or VM remains a far capability.  Its \kw{actorStackQuery} result
is transported by the remote-representation judgment; object references in a
frame become values, far references, or authorized mirrors rather than raw
near references.  A remote controller first applies
Equation~\eqref{eq:pause-actor-effect}, reads the resulting paused stack
snapshot, and then names its pause token in
Equation~\eqref{eq:replace-paused-actor-stack},
\eqref{eq:resume-paused-actor-full-speed}, or
\eqref{eq:replace-resume-paused-actor-stack}.  Formally, a remote debugger step
has the component-wise shape
\begin{equation}\label{eq:remote-debugger-component-step}
 \frac{\mathfrak F(w),A,g\vdash\mathsf{debugRequest}(\delta)
          \Downarrow\mathfrak F'(w)}
      {\mathfrak F\longmapsto_{\mathsf{remoteDebug}}
       \mathfrak F[w\mapsto\mathfrak F'(w)]},
 \qquad
 \forall v\ne w.\ \mathfrak F'(v)=\mathfrak F(v).
\end{equation}
The request is delivered using the ordinary cross-VM transport and executes
the corresponding local reflective command at $w$'s safepoint.  There is no
distributed atomic transaction spanning several VMs: coordinated multi-actor
debugging is a protocol of individually atomic pause, inspect, install, and
resume operations.  Pause tokens reject delayed or replayed commands after
either an intervening edit or a resume.

Finally, VM primitives are ordinary messages to a mirror whose record has
target $\kw{vmTarget}(w)$ and right $\kw{primitive}$.  Their platform-specific state
transformers lie outside this language semantics, but the capability premise
is mandatory:
\[
\dfrac{\begin{gathered}
       \widehat H_A(g)=\kw{mirror}\langle C_g,w,\kw{vmTarget}(w),R\rangle
       \qquad\kw{primitive}\in R\\
       \mathsf{hostPrimitive}(\Omega,s,\bar o)=\langle\Omega',v\rangle
       \end{gathered}}
 {\Omega,A,g\vdash\mathsf{invokeVMPrimitive}(s,\bar o)
       \Downarrow\langle\Omega',v\rangle}
\quad(\rulelabel{VM-Primitive})
\]
There is no primitive-call syntax and no bypass around the mirror record.
